% main.tex
% Visual Analytics for Multi-Criteria EV Charging Infrastructure Siting
% and Scenario Exploration

\documentclass[journal]{IEEEtran}

\usepackage[utf8]{inputenc}
\usepackage[T1]{fontenc}
\usepackage{amsmath,amssymb}
\usepackage{graphicx}
\usepackage{booktabs}
\usepackage{enumitem}
\usepackage{xcolor}
\usepackage{hyperref}
\usepackage{url}
\usepackage{subcaption}
\usepackage{verbatim}
\usepackage{pifont}
\usepackage{multirow}

\hypersetup{
  colorlinks=true,
  linkcolor=black,
  citecolor=blue!70!black,
  urlcolor=blue!70!black
}

\title{Interactive Visual Analytics for Multi-Criteria EV Charging Infrastructure Siting and Scenario Exploration}

\author{
  \IEEEauthorblockN{Author Names}
  \IEEEauthorblockA{Affiliation \\
  Email}
}

\begin{document}

\maketitle

\begin{abstract}
Planning urban infrastructure often requires balancing competing objectives such as accessibility, demand, grid capacity, and social equity. Multi-criteria decision analysis (MCDA) is frequently used to support such planning tasks, yet traditional workflows remain largely static and opaque, limiting stakeholders' ability to explore trade-offs and understand how modelling assumptions influence outcomes. We present a visual analytics platform that integrates interactive multi-criteria suitability analysis with real-time what-if scenario evaluation for electric vehicle charging point (EVCP) siting. The system combines coordinated visual views---including parallel coordinate weight controls, hexagonal spatial suitability maps with multi-variate glyph overlays, and diagnostic panels---enabling planners to interactively adjust criterion weights (via direct manipulation or Analytic Hierarchy Process pairwise comparison), examine ranking sensitivity, and explore the spatial implications of alternative planning priorities. Beyond site selection, the platform incorporates a lightweight impact estimation module that computes energy delivery, carbon reduction, grid headroom, and socio-demographic coverage for proposed charger deployments. These impact estimates are overlaid on Distribution Future Energy Scenarios (DFES) timeseries projections, transforming passive reference data into temporal stress tests that reveal when proposed infrastructure becomes insufficient under different demand pathways. Users can save, compare, and export complete planning scenarios via a radar chart comparison view. The entire analytical pipeline---including MCDA scoring via Weighted Sum, Weighted Product, and TOPSIS methods---runs client-side in DuckDB-WASM over approximately 118,000 H3 hexagonal cells, eliminating server infrastructure requirements while maintaining interactive response times. We demonstrate the approach through a case study of charging infrastructure planning in Greater London using heterogeneous open datasets. The results illustrate how interactive visual analytics can support transparent exploration of planning trade-offs and facilitate more informed infrastructure decision-making.
\end{abstract}

\begin{IEEEkeywords}
Visual analytics, multi-criteria decision analysis, electric vehicle charging infrastructure, spatial decision support, what-if analysis, DuckDB, H3 hexagonal index
\end{IEEEkeywords}

% 1_introduction.tex

\section{Introduction}
\label{sec:introduction}

The transition to electric mobility is a central pillar of urban decarbonisation strategies worldwide. In the United Kingdom, the government has committed to ending the sale of new petrol and diesel vehicles by 2035, placing significant pressure on local authorities and network operators to expand public electric vehicle charging point (EVCP) infrastructure at pace~\cite{tfl2025evstrategy}. London alone targets over 40,000 public charge points by 2030, yet as of early 2025 only a fraction of that target has been realised, with substantial geographic disparities between boroughs. The siting of new EVCP infrastructure is a complex planning problem: candidate locations must balance competing criteria including population density, existing charging coverage, electricity grid capacity, traffic demand, accessibility to services, and socio-economic equity~\cite{banegas2023systematic,kaya2020evcs}.

Multi-criteria decision analysis (MCDA) provides a principled framework for evaluating alternatives against multiple, often conflicting, objectives~\cite{malczewski2015gismcda}. When combined with geographic information systems (GIS), MCDA methods such as the Weighted Sum Model (WSM), Weighted Product Model (WPM), and TOPSIS~\cite{hwang1981topsis} enable spatially explicit suitability assessments that rank candidate locations based on user-defined priorities~\cite{malczewski1999gis}. The Analytic Hierarchy Process (AHP)~\cite{saaty1980ahp} offers a structured elicitation procedure for deriving criterion weights from pairwise comparisons, improving the transparency and consistency of weight assignment.

Despite the maturity of these analytical methods, their application in infrastructure planning typically follows a static, batch-oriented workflow: analysts define weights, execute a GIS model, inspect the output, and iterate by adjusting parameters and re-running the pipeline. This process is opaque to non-technical stakeholders, obscures the sensitivity of results to weight choices, and provides no mechanism for evaluating the downstream consequences of siting decisions---such as the energy delivered, carbon saved, or grid stress induced by a proposed deployment~\cite{andrienko2007geovisual,jankowski2001gis}.

Visual analytics offers an alternative paradigm that combines computational analysis with interactive visual representations, enabling analysts to steer analytical processes, explore parameter spaces, and develop insight through direct engagement with data~\cite{keim2008visual}. The visual analytics mantra---\emph{analyse first, show the important, zoom, filter, analyse further, details on demand}---provides a design framework well suited to the iterative, exploratory nature of infrastructure planning~\cite{shneiderman1996eyes}. However, existing visual analytics systems for spatial MCDA remain limited in two respects. First, most focus exclusively on the site selection phase and do not integrate impact estimation or scenario simulation. Second, few operate entirely client-side, instead relying on server infrastructure that constrains deployment flexibility and introduces latency in interactive weight exploration.

In this paper, we present a visual analytics platform that addresses both limitations. The system integrates three analytical stages within a coordinated multiple-view interface:

\begin{enumerate}[nosep]
    \item \textbf{Multi-criteria suitability analysis}: interactive weight adjustment (via direct manipulation or AHP pairwise comparison) with real-time MCDA scoring across approximately 118,000 H3 hexagonal cells using three complementary methods (WSM, WPM, TOPSIS).

    \item \textbf{What-if impact simulation}: a lightweight analytical model that estimates energy delivery, carbon reduction, grid headroom impact, population coverage, and financial indicators for proposed charger deployments at user-selected locations.

    \item \textbf{Temporal contextualisation}: overlay of scenario-level impact estimates on UK Power Networks' Distribution Future Energy Scenarios (DFES) demand projections, revealing the year at which proposed infrastructure becomes insufficient under different demand pathways.
\end{enumerate}

The entire analytical pipeline runs client-side in a DuckDB-WASM instance, processing Parquet-encoded spatial data without server infrastructure. This architecture enables interactive response times during weight exploration while supporting deployment as a static web application.

We demonstrate the platform through a case study of EVCP siting in Greater London, using nine heterogeneous open datasets spanning census demographics, transport emissions, electricity network capacity, accessibility indices, and existing charger locations. The case study illustrates a complete planning workflow: a planner adjusts criterion weights, explores suitability patterns on the map via choropleth and multi-variate glyph encodings, places candidate chargers, evaluates their impact, examines how the deployment performs against future demand trajectories, and compares alternative scenarios.

The contributions of this work are:

\begin{itemize}[nosep]
    \item A visual analytics design pattern that integrates MCDA site selection, what-if impact simulation, and temporal scenario contextualisation within a coordinated multiple-view architecture.
    \item A client-side analytical pipeline using DuckDB-WASM and H3 hexagonal indexing that enables real-time MCDA computation over large spatial datasets without server infrastructure.
    \item A demonstration of the approach through a real-world EVCP planning case study in Greater London, grounded in heterogeneous open datasets.
\end{itemize}

% 2_related_work.tex

\section{Related Work}
\label{sec:related-work}

This work draws on three research areas: visual analytics for spatial decision support, multi-criteria decision analysis in GIS, and methodologies for electric vehicle charging infrastructure siting. We review each area and position our contribution relative to existing approaches.

%--------------------------------------------------------------
\subsection{Visual Analytics for Spatial Decision Support}
\label{sec:rw-va}

Visual analytics combines automated analysis with interactive visual representations to support analytical reasoning over complex data~\cite{keim2008visual}. In the spatial domain, Andrienko et al.~\cite{andrienko2007geovisual} established a research agenda for geovisual analytics that emphasises the synergy between computational techniques and human analytical capabilities, arguing that interactive exploration of spatial patterns, relationships, and trends is essential for informed decision-making.

Coordinated multiple views (CMV) are a foundational design pattern in visual analytics, enabling users to examine data through multiple linked representations~\cite{roberts2007coordinated}. In spatial decision support, CMV architectures typically coordinate a map view with statistical summaries, parameter controls, and diagnostic displays, allowing analysts to observe how changes in one dimension propagate across others. Shneiderman's information-seeking mantra---\emph{overview first, zoom and filter, details on demand}~\cite{shneiderman1996eyes}---provides a complementary interaction framework that guides the progressive disclosure of spatial information.

Several systems have applied visual analytics principles to spatial planning problems. Jankowski and Nyerges~\cite{jankowski2001gis} developed GIS-based collaborative decision support tools that combine spatial data with group deliberation processes. Rinner and Malczewski~\cite{rinner2002owa} demonstrated web-enabled spatial decision analysis using Ordered Weighted Averaging operators with interactive map visualisation. More recent work has explored visual analytics for scenario planning~\cite{andrienko2003exploratory} and what-if analysis in infrastructure contexts, though few systems integrate site selection with downstream impact estimation in a single interactive environment.

Our platform extends this line of work by coupling multi-criteria spatial analysis with a what-if simulation layer and temporal scenario contextualisation, all within a coordinated multiple-view interface that operates entirely client-side.

%--------------------------------------------------------------
\subsection{GIS-Based Multi-Criteria Decision Analysis}
\label{sec:rw-mcda}

The integration of MCDA with GIS has been extensively studied over the past three decades. Malczewski's foundational work~\cite{malczewski1999gis} established the conceptual framework for combining spatial data with multi-criteria evaluation methods, while Malczewski and Rinner~\cite{malczewski2015gismcda} provided a comprehensive treatment of GIS-MCDA at an advanced level, covering decision models, value scaling, criterion weighting, and combination rules.

Among weighting methods, the Analytic Hierarchy Process (AHP)~\cite{saaty1980ahp} remains widely used for its structured elicitation of preferences through pairwise comparison. AHP produces a consistent weight vector from a reciprocal comparison matrix and provides a built-in consistency check (the Consistency Ratio), making it suitable for participatory settings where decision-makers need guidance on the coherence of their judgements. Several GIS-MCDA systems have incorporated AHP for weight derivation, though typically as an offline preprocessing step rather than as an interactive, integrated component~\cite{priefer2022evcs,kaya2020evcs}.

Aggregation methods vary in their treatment of criteria trade-offs. The Weighted Sum Model (WSM) is the most common, offering full compensability: a high score on one criterion can compensate for a low score on another. The Weighted Product Model (WPM) uses multiplicative aggregation, which penalises alternatives with very low values on any criterion~\cite{malczewski2015gismcda}. TOPSIS~\cite{hwang1981topsis} ranks alternatives by their simultaneous distance to an ideal best and ideal worst solution, providing a geometric perspective on trade-off structure. Offering all three methods within a single interactive environment enables analysts to assess the robustness of rankings to methodological choice---a form of sensitivity analysis rarely available in batch-oriented GIS-MCDA workflows.

Our system integrates AHP weight elicitation as an interactive component within the dashboard, enabling real-time pairwise comparison alongside direct weight manipulation, and supports all three aggregation methods with immediate visual feedback.

%--------------------------------------------------------------
\subsection{Electric Vehicle Charging Infrastructure Siting}
\label{sec:rw-evcp}

The siting of electric vehicle charging stations has attracted growing research attention as nations accelerate the transition to electric mobility. Banegas and Mamkhezri~\cite{banegas2023systematic} conducted a systematic review of GIS-based methods and criteria for EVCS site selection, examining 52 studies published between 2010 and 2021. They identified geographic and demographic variables as the most frequently used factors and noted that map algebra and overlay methods dominate the analytical approaches. Their review categorised siting criteria into six groups: geographic, energy, charging station attributes, EV characteristics, transportation and traffic, economic, socio-demographic, and environmental factors.

Kaya et al.~\cite{kaya2020evcs} applied AHP combined with PROMETHEE and VIKOR methods using 19 parameters organised into economic, geographical, energy, social, and transportation categories. Priefer and Steiger~\cite{priefer2022evcs} used AHP with GIS overlay for suitability analysis, demonstrating how pairwise comparison can factorise weight assignments for EVCS parameters including accessibility, proximity to users, environmental impact, and power grid capacity.

A consistent finding across this literature is that EVCS siting is inherently multi-criteria and context-dependent: the relative importance of demand indicators, infrastructure constraints, equity considerations, and environmental factors varies by planning context and stakeholder priorities~\cite{banegas2023systematic}. This observation motivates the interactive weight exploration approach adopted in our platform, where planners can adjust criterion priorities and immediately observe the spatial consequences.

However, existing EVCS siting studies focus almost exclusively on identifying suitable locations. They do not evaluate the downstream impacts of proposed deployments---such as energy delivered, carbon reduction potential, or grid stress---nor do they contextualise siting decisions against forward-looking demand projections. Our work addresses this gap by integrating site selection with a lightweight impact estimation model and temporal scenario analysis through Distribution Future Energy Scenarios (DFES).

% 3_methodology.tex

\section{Analytical Framework}
\label{sec:methodology}

The platform integrates two analytical components: a multi-criteria suitability model for ranking candidate locations and a lightweight impact estimation model for evaluating the consequences of proposed charger deployments. Both operate client-side over a hexagonal spatial tessellation using H3 cells at resolution~10, with each cell covering approximately $0.015~\text{km}^2$.

%--------------------------------------------------------------
\subsection{Spatial Data Model}
\label{sec:spatial-data}

The study area (Greater London) is discretised into $n \approx 118{,}000$ H3 hexagonal cells at resolution~10 using Uber's H3 hierarchical spatial index~\cite{uber2018h3}. Each cell $h$ carries a vector of $m = 9$ criterion values derived from heterogeneous urban datasets:

\begin{equation}
\mathbf{x}_h = \bigl(x_{h,1},\; x_{h,2},\; \ldots,\; x_{h,m}\bigr)
\end{equation}

\noindent All criterion values are pre-normalised to $[0, 1]$ using min-max normalisation prior to ingestion, ensuring deterministic and reproducible scoring without runtime normalisation cost. Table~\ref{tab:criteria} lists the nine criteria and their data sources. Each cell also carries three administrative geography attributes---LSOA code, LSOA name, and London borough name---enabling direct linkage to external datasets such as DFES projections (Section~\ref{sec:dfes}) without runtime spatial joins.

\begin{table}[htbp]
\centering
\caption{MCDA criteria used in the suitability model.}
\label{tab:criteria}
\small
\begin{tabular}{clll}
\toprule
\textbf{\#} & \textbf{Criterion} & \textbf{Category} & \textbf{Source} \\
\midrule
$C_1$ & Population density & Demand & ONS Census 2021~\cite{ons2021census} \\
$C_2$ & Car ownership ($\geq 1$ vehicle) & Demand & ONS Census 2021~\cite{ons2021census} \\
$C_3$ & Deprivation ($\geq 2$ dimensions) & Equity & ONS Census 2021~\cite{ons2021census} \\
$C_4$ & Employment accessibility (30\,min PT) & Access & Verduzco Torres et al.~\cite{verduzco2024accessibility} \\
$C_5$ & Supermarket accessibility (30\,min PT) & Access & Verduzco Torres et al.~\cite{verduzco2024accessibility} \\
$C_6$ & Road transport CO\textsubscript{2} emissions & Environment & LAEI 2022~\cite{laei2022emissions} \\
$C_7$ & Grid demand headroom & Infrastructure & UKPN substation data \\
$C_8$ & Motorised traffic index & Demand & DfT traffic counts \\
$C_9$ & Distance to nearest existing EVCP & Coverage & Open Charge Map~\cite{openchargemapdata} \\
\bottomrule
\end{tabular}
\end{table}

The hexagonal tessellation was chosen over regular square grids for several reasons. Hexagons minimise edge effects in spatial aggregation because each cell has six equidistant neighbours, reducing directional bias in adjacency-based analyses~\cite{uber2018h3}. The H3 system provides a hierarchical structure: cells at resolution~10 (the base analytical unit) can be aggregated to coarser parent resolutions ($r \in \{5, 6, 7, 8, 9\}$) via the \texttt{cellToParent} function, enabling multi-scale visualisation without recomputing MCDA scores.

%--------------------------------------------------------------
\subsection{Multi-Criteria Suitability Model}
\label{sec:mcda-model}

The platform supports three complementary MCDA scoring methods, each offering a different perspective on the trade-off structure. In all methods, the user assigns a weight vector $\mathbf{w} = (w_1, w_2, \ldots, w_m)$ subject to the constraint $\sum_{i=1}^{m} w_i = 1$, with $w_i \geq 0$. Each criterion can be designated as a \emph{benefit} criterion (higher values are preferred) or a \emph{cost} criterion (lower values are preferred); cost criteria are inverted ($1 - \bar{x}_{h,i}$) before scoring.

%--- WSM ---
\subsubsection{Weighted Sum Model (WSM)}

The Weighted Sum Model computes a suitability score as a linear combination of normalised criterion values:

\begin{equation}
\label{eq:wsm}
S_h^{\text{WSM}} = \sum_{i=1}^{m} w_i \cdot \bar{x}_{h,i}
\end{equation}

\noindent where $\bar{x}_{h,i} \in [0,1]$ is the normalised value of criterion~$i$ at cell~$h$. The WSM provides full compensability: a high score on one criterion can offset a low score on another. The score range is $[0, 1]$ when all normalised values and weights are non-negative.

%--- WPM ---
\subsubsection{Weighted Product Model (WPM)}

The Weighted Product Model uses multiplicative aggregation, which reduces compensation between criteria:

\begin{equation}
\label{eq:wpm}
S_h^{\text{WPM}} = \prod_{i=1}^{m} \bigl(\bar{x}_{h,i}\bigr)^{w_i}
\end{equation}

\noindent A small floor value ($\epsilon = 0.001$) is applied to prevent zero scores from nullifying the product: $\bar{x}_{h,i}' = \max(\bar{x}_{h,i},\; \epsilon)$. The WPM penalises cells with extremely low values on any criterion more heavily than the WSM, producing rankings that are more sensitive to worst-case performance.

%--- TOPSIS ---
\subsubsection{TOPSIS}

The Technique for Order Preference by Similarity to Ideal Solution (TOPSIS)~\cite{hwang1981topsis} ranks alternatives based on their geometric distance to idealised best and worst solutions.

First, a weighted normalised matrix is constructed:

\begin{equation}
v_{h,i} = w_i \cdot \bar{x}_{h,i}
\end{equation}

\noindent The positive ideal solution (PIS) and negative ideal solution (NIS) are determined:

\begin{equation}
A^{+}_i = \max_{h} \; v_{h,i}, \qquad
A^{-}_i = \min_{h} \; v_{h,i}
\end{equation}

\noindent Euclidean distances from each cell to the PIS and NIS are:

\begin{equation}
D_h^{+} = \sqrt{\sum_{i=1}^{m} (v_{h,i} - A^{+}_i)^2}, \qquad
D_h^{-} = \sqrt{\sum_{i=1}^{m} (v_{h,i} - A^{-}_i)^2}
\end{equation}

\noindent The relative closeness to the ideal solution yields the TOPSIS score:

\begin{equation}
\label{eq:topsis}
S_h^{\text{TOPSIS}} = \frac{D_h^{-}}{D_h^{+} + D_h^{-}}
\end{equation}

\noindent with $S_h^{\text{TOPSIS}} \in [0,1]$. A cell scoring close to 1 is nearest to the ideal best and farthest from the ideal worst.

%--------------------------------------------------------------
\subsection{Weight Elicitation via AHP}
\label{sec:ahp}

Criterion weights can be derived interactively using the Analytic Hierarchy Process (AHP)~\cite{saaty1980ahp}. Users perform pairwise comparisons of criteria importance, populating an $m \times m$ reciprocal matrix $\mathbf{A}$ where:

\begin{equation}
a_{ij} = \frac{1}{a_{ji}}, \qquad a_{ii} = 1
\end{equation}

\noindent Each entry $a_{ij}$ represents the relative importance of criterion~$i$ over criterion~$j$ on the Saaty fundamental scale ($\tfrac{1}{9}$ to $9$).

\paragraph{Weight derivation.}
Priority weights are derived using the Row Geometric Mean Method (RGMM):

\begin{equation}
\label{eq:ahp-weights}
\hat{w}_i = \frac{\left(\prod_{j=1}^{m} a_{ij}\right)^{1/m}}{\sum_{k=1}^{m} \left(\prod_{j=1}^{m} a_{kj}\right)^{1/m}}
\end{equation}

\paragraph{Consistency checking.}
The consistency of the pairwise comparison matrix is assessed through the principal eigenvalue $\lambda_{\max}$:

\begin{equation}
\lambda_{\max} = \frac{1}{m} \sum_{i=1}^{m} \frac{(\mathbf{A} \hat{\mathbf{w}})_i}{\hat{w}_i}
\end{equation}

\noindent The Consistency Index (CI) and Consistency Ratio (CR) are:

\begin{equation}
\text{CI} = \frac{\lambda_{\max} - m}{m - 1}, \qquad
\text{CR} = \frac{\text{CI}}{\text{RI}(m)}
\end{equation}

\noindent where $\text{RI}(m)$ is the Random Index for a matrix of size~$m$. A comparison matrix is considered acceptably consistent if $\text{CR} < 0.10$.

%--------------------------------------------------------------
\subsection{In-Browser Analytical Pipeline}
\label{sec:data-pipeline}

All spatial data processing runs client-side in a DuckDB-WASM instance~\cite{duckdb2019,duckdbwasm2022}, eliminating server infrastructure requirements. Ten Parquet files (one per criterion layer plus supplementary layers) are fetched on page load, registered as DuckDB tables, and joined on the H3 cell identifier into a unified \texttt{mcda\_base} view:

\begin{equation}
\texttt{mcda\_base} = \underset{i=1}{\overset{10}{\bowtie}} \; T_i \quad \text{on } \texttt{h3\_cell}
\end{equation}

\noindent Each Parquet file carries administrative geography fields (\texttt{lsoa21cd}, \texttt{lsoa21nm}, \texttt{borough\_name}) that are exposed once from the base table. The MCDA scoring query is dynamically generated from the current weight vector and selected method, executed against the \texttt{mcda\_base} view, and returns both the composite score and raw criterion values per cell in a single pass.

This architecture offers several advantages. Parquet's columnar storage enables efficient scans over the subset of columns required by the active criteria. DuckDB's vectorised execution engine processes the approximately 118,000-row scoring query in under 200\,ms on modern hardware, well within the interactive latency threshold. The H3 spatial extension loaded into DuckDB provides \texttt{h3\_cell\_to\_parent} for resolution rollup queries, enabling multi-scale aggregation directly within SQL.

%--------------------------------------------------------------
\subsection{Hierarchical Spatial Aggregation}
\label{sec:aggregation}

For interactive visualisation at varying zoom levels, MCDA scores computed at the base resolution (H3 resolution~10) are aggregated to coarser parent resolutions ($r \in \{5, 6, 7, 8, 9\}$) using mean pooling:

\begin{equation}
S_{\text{parent}}^{(r)} = \frac{1}{|\mathcal{C}_{\text{parent}}|} \sum_{h \in \mathcal{C}_{\text{parent}}} S_h
\end{equation}

\noindent where $\mathcal{C}_{\text{parent}}$ is the set of resolution-10 child cells belonging to the parent cell at resolution~$r$. The display resolution is selected dynamically based on the current map zoom level (e.g., zoom~14+ maps to resolution~10, zoom~8 to resolution~7), reducing the number of rendered hexagons from $\sim\!118{,}000$ at full resolution to $\sim\!2{,}000$ at typical overview zoom levels.


%==============================================================
\section{Impact Estimation Model}
\label{sec:impact}

Once a user selects a candidate H3 cell and configures a charger deployment (type and count), the platform estimates the potential impacts of that intervention using a set of lightweight analytical equations. These are designed for real-time evaluation within an interactive environment, prioritising tractability and transparency over simulation fidelity.

%--------------------------------------------------------------
\subsection{Charger Specifications}
\label{sec:charger-specs}

The platform supports four charger types reflecting the current UK public charging market:

\begin{table}[htbp]
\centering
\caption{Charger type specifications.}
\label{tab:charger-specs}
\small
\begin{tabular}{llrrl}
\toprule
\textbf{Type} & \textbf{Label} & \textbf{Power (kW)} & \textbf{Cost (\pounds/unit)} & \textbf{Install} \\
\midrule
Slow & AC Slow & 3.7 & 1{,}000 & 1 month \\
Fast & AC Fast & 22 & 5{,}000 & 2 months \\
Rapid & DC Rapid & 50 & 40{,}000 & 4 months \\
Ultra-rapid & DC Ultra-rapid & 150 & 100{,}000 & 6 months \\
\bottomrule
\end{tabular}
\end{table}

%--------------------------------------------------------------
\subsection{Demand and Impact Estimation}
\label{sec:demand}

Local daily charging demand is estimated from the spatial attributes of the selected cell:

\begin{equation}
\label{eq:demand}
D_h = \rho_h \cdot A_{\text{cell}} \cdot \gamma_h \cdot \alpha_{\text{car}} \cdot \alpha_{\text{EV}} \cdot \alpha_{\text{charge}} \cdot E_{\text{session}}
\end{equation}

\noindent where $\rho_h$ is population density, $A_{\text{cell}} = 0.015~\text{km}^2$, $\gamma_h$ is the local car ownership rate, $\alpha_{\text{car}} = 0.4$ is the cars-to-population ratio, $\alpha_{\text{EV}} = 0.15$ is the EV adoption rate, $\alpha_{\text{charge}} = 0.3$ is the daily public charging probability, and $E_{\text{session}} = 30~\text{kWh}$ is the average energy per session.

The utilisation factor captures the balance between local demand and installed capacity:

\begin{equation}
\label{eq:utilisation}
U = \min\!\left(1,\; \frac{D_h}{N \cdot P \cdot H}\right)
\end{equation}

\noindent where $N$ is the number of chargers, $P$ is the rated power (kW), and $H = 16$~hours is the daily operating window. Annual energy delivered is $E_{\text{annual}} = N \cdot P \cdot H \cdot U \cdot 365$ (kWh/year).

Carbon savings compare displaced ICE emissions with grid emissions from EV charging:

\begin{equation}
\label{eq:carbon}
C_{\text{saved}} = \frac{E_{\text{annual}}}{\eta_{\text{EV}}} \cdot \text{EF}_{\text{ICE}} - E_{\text{annual}} \cdot \text{EF}_{\text{grid}}
\end{equation}

\noindent where $\eta_{\text{EV}} = 0.18~\text{kWh/km}$, $\text{EF}_{\text{ICE}} = 0.21~\text{kgCO}_2\text{/km}$, and $\text{EF}_{\text{grid}} = 0.233~\text{kgCO}_2\text{/kWh}$.

Grid impact is estimated as the peak demand ($P_{\text{peak}} = N \cdot P \cdot f_{\text{div}}$, with diversity factor $f_{\text{div}} = 0.7$) expressed as a fraction of available substation capacity. Population served is estimated over the target cell and its six H3 neighbours (the 1-ring). Financial indicators include total installation cost and estimated annual revenue at a public charging tariff of \pounds0.35/kWh.

Table~\ref{tab:kpis} summarises the eight key performance indicators computed for each scenario.

\begin{table}[htbp]
\centering
\caption{Impact KPIs reported by the platform.}
\label{tab:kpis}
\small
\begin{tabular}{llll}
\toprule
\textbf{KPI} & \textbf{Unit} & \textbf{Equation} & \textbf{Aggregation} \\
\midrule
Energy delivered & kWh/year & Eq.~\ref{eq:demand}--\ref{eq:utilisation} & Sum \\
Carbon saved & tCO\textsubscript{2}/year & Eq.~\ref{eq:carbon} & Sum \\
Installation cost & \pounds & per-unit cost & Sum \\
Annual revenue & \pounds/year & energy $\times$ tariff & Sum \\
Peak demand & kW & $N \cdot P \cdot f_{\text{div}}$ & Sum \\
Headroom impact & \% & peak / capacity & Average \\
Population served & persons & density $\times$ area $\times$ 7 & Sum \\
Utilisation factor & \% & Eq.~\ref{eq:utilisation} & Average \\
\bottomrule
\end{tabular}
\end{table}

When the user places chargers at multiple locations, additive metrics are summed across placements and ratio metrics are averaged.

% 4_system_design.tex

\section{Visual Analytics Design}
\label{sec:vis-design}

The dashboard follows a coordinated multiple view (CMV) architecture~\cite{roberts2007coordinated} in which spatial, analytical, and temporal perspectives on EVCP siting are tightly linked through shared state. The design is guided by the visual analytics mantra~\cite{keim2008visual}---\emph{analyse first, show the important, zoom, filter, analyse further, details on demand}---and draws on established principles for multi-criteria spatial decision support~\cite{andrienko2007geovisual,jankowski2001gis}.

%--------------------------------------------------------------
\subsection{Dashboard Layout}
\label{sec:layout}

The interface is organised into three coordinated panels (Figure~\ref{fig:dashboard-layout}):

\begin{enumerate}[nosep]
    \item \textbf{Weight Laboratory} (left panel): interactive controls for MCDA method selection, criterion weighting via parallel coordinates, pairwise AHP comparison, and a weight diagnostics view.
    \item \textbf{Spatial View} (centre panel): a web map displaying H3 hexagonal choropleth and glyph overlays, with a DFES timeseries panel anchored below.
    \item \textbf{Impact Panel} (right panel): scenario placement list, eight KPI cards, a radar chart for scenario comparison, and a per-LSOA breakdown of aggregated impacts.
\end{enumerate}

\noindent All panels share state through a reactive store architecture (Zustand), ensuring that weight changes, map interactions, and scenario edits propagate immediately across views without explicit event wiring.

\begin{figure}[htbp]
    \centering
    % PLACEHOLDER: Screenshot of the full dashboard showing all three panels
    \fbox{\parbox{0.95\columnwidth}{\centering\vspace{3cm}\textbf{[Figure: Dashboard Layout]}\\\small Full dashboard view showing the Weight Laboratory (left), Spatial View with H3 choropleth (centre), and Impact Panel with KPI cards (right).\vspace{0.5cm}}}
    \caption{The dashboard layout with three coordinated panels. The Weight Laboratory provides interactive weight controls and AHP pairwise comparison. The Spatial View renders MCDA suitability scores as a hexagonal choropleth with optional glyph overlays. The Impact Panel displays KPIs and scenario comparison charts.}
    \label{fig:dashboard-layout}
\end{figure}

%--------------------------------------------------------------
\subsection{Weight Laboratory}
\label{sec:weight-lab}

The Weight Laboratory provides two complementary mechanisms for criterion weight assignment:

\paragraph{Direct manipulation via parallel coordinates.}
A parallel coordinate display renders the $m$ active criteria as vertical axes. Each axis carries a draggable handle that controls the criterion weight; dragging upward increases the weight while the remaining weights are proportionally adjusted to maintain the unit-sum constraint. A horizontal bar chart below the parallel coordinates provides a secondary encoding of the weight distribution. Criteria can be toggled on or off, and their benefit/cost polarity can be switched, with each change triggering an immediate MCDA re-computation.

\paragraph{Pairwise comparison via AHP.}
An interactive AHP matrix interface (Figure~\ref{fig:ahp-interface}) enables structured weight elicitation. Users adjust pairwise importance ratios on the Saaty scale ($\tfrac{1}{9}$ to $9$) using sliders. The system solves the comparison matrix in real time using the Row Geometric Mean Method (Section~\ref{sec:ahp}), displaying the derived weights, the Consistency Ratio (CR), and a consistency status indicator. When $\text{CR} < 0.10$, the derived weights can be applied to the MCDA model with a single action, replacing the current direct-manipulation weights. A diagnostics panel shows the full comparison matrix, enabling users to identify and correct inconsistent judgements.

\begin{figure}[htbp]
    \centering
    % PLACEHOLDER: Screenshot of the AHP comparison interface
    \fbox{\parbox{0.95\columnwidth}{\centering\vspace{3cm}\textbf{[Figure: AHP Interface]}\\\small The AHP pairwise comparison panel showing sliders for each criterion pair, derived weights, and the Consistency Ratio indicator.\vspace{0.5cm}}}
    \caption{The AHP pairwise comparison interface. Sliders control the relative importance of each criterion pair. The derived weight vector and Consistency Ratio are computed and displayed in real time.}
    \label{fig:ahp-interface}
\end{figure}

\paragraph{Method selection.}
A dropdown control allows switching between WSM, WPM, and TOPSIS aggregation methods. Changing the method triggers a new DuckDB query with the same weight vector but a different scoring function, enabling rapid comparison of how methodological choice affects the spatial ranking pattern.

%--------------------------------------------------------------
\subsection{Spatial View and Glyph Encoding}
\label{sec:spatial-view}

The primary spatial encoding is a hexagonal choropleth rendered as GeoJSON polygons over a CartoDB Positron basemap using MapLibre GL JS~\cite{maplibregl}. Each H3 cell is coloured by its MCDA suitability score using a perceptually uniform \texttt{viridis} colour scale, with opacity adjustable via a continuous slider.

An optional glyph overlay (Figure~\ref{fig:glyph-map}) projects per-cell criterion profiles as either \emph{bar glyphs} (up to 6 criteria) or \emph{rose glyphs} (up to 8 criteria), drawn on HTML5 Canvas and positioned at cell centroids. The glyph layer encodes multiple visual variables following Bertin's~\cite{bertin1983semiology} framework:

\begin{itemize}[nosep]
    \item \textbf{Size}: petal length (rose) or bar height encodes the normalised criterion value.
    \item \textbf{Width}: petal angular sweep (rose only) encodes the criterion weight, making higher-weighted criteria visually dominant.
    \item \textbf{Texture}: a diagonal hatch pattern distinguishes \emph{cost} criteria from \emph{benefit} criteria within the same glyph, following the convention that texture signals qualitative difference.
    \item \textbf{Comparison mode}: clicking a glyph pins it as a reference; subsequent hover glyphs show a dashed outline of the pinned profile for direct pairwise comparison.
\end{itemize}

\noindent The glyph aggregation resolution is user-controlled (H3~r6--r10), enabling analysts to switch between local detail and regional overview without recomputing MCDA scores.

\begin{figure}[htbp]
    \centering
    % PLACEHOLDER: Screenshot showing glyph overlay on map
    \fbox{\parbox{0.95\columnwidth}{\centering\vspace{3cm}\textbf{[Figure: Glyph Map]}\\\small Map view with rose glyph overlay showing per-cell criterion profiles. Cost criteria are distinguished by a hatch pattern. A pinned reference glyph (solid) enables comparison with a hovered glyph (dashed outline).\vspace{0.5cm}}}
    \caption{The glyph overlay showing rose glyphs at H3 resolution~8. Each petal represents a criterion: length encodes the normalised value, angular width encodes the weight, and hatching distinguishes cost criteria. The comparison mode overlays a pinned reference profile (solid) against the hovered cell (dashed).}
    \label{fig:glyph-map}
\end{figure}

%--------------------------------------------------------------
\subsection{Simulation Mode and Placement Interaction}
\label{sec:simulation-mode}

Activating \emph{Simulation Mode} changes the map cursor to an EVCP placement marker and enables a click-to-place interaction on any H3 cell. On click, the system:

\begin{enumerate}[nosep]
    \item Retrieves the MCDA query result row for the clicked cell, extracting both raw and normalised criterion values as well as LSOA and borough metadata from the underlying DuckDB tables.
    \item Opens a placement configuration popover displaying the reverse-geocoded address (via Nominatim), a charger type selector, a charger count slider, and the estimated installation cost.
    \item On confirmation, stores the placement with its full spatial context---criterion values, LSOA code, LSOA name, and borough name---ensuring that impact calculations remain reproducible even if the user later changes MCDA weights.
\end{enumerate}

\noindent This \emph{capture-at-placement-time} design is a deliberate architectural choice: by snapshotting the spatial data when the placement is created, the impact model is decoupled from the current MCDA state, preventing silent invalidation of impact estimates when criterion weights change.

%--------------------------------------------------------------
\subsection{Impact Panel: KPIs and Radar Comparison}
\label{sec:impact-panel}

The right panel (Figure~\ref{fig:impact-panel}) renders eight KPI cards (Section~\ref{sec:demand}) with large-format numbers and contextual labels. When multiple scenarios are saved, a radar chart normalises six axes---energy, carbon, revenue, population, utilisation, and headroom spare---to $[0, 1]$ using the maximum value across all series, enabling direct visual comparison of scenario profiles.

\begin{figure}[htbp]
    \centering
    % PLACEHOLDER: Screenshot of Impact Panel with KPI cards and radar chart
    \fbox{\parbox{0.95\columnwidth}{\centering\vspace{3cm}\textbf{[Figure: Impact Panel]}\\\small The Impact Panel showing KPI cards (top), a radar chart comparing two saved scenarios (middle), and a per-LSOA breakdown table (bottom).\vspace{0.5cm}}}
    \caption{The Impact Panel. KPI cards display energy delivered, carbon saved, installation cost, revenue, peak demand, headroom impact, population served, and utilisation. The radar chart enables multi-dimensional comparison of saved scenarios.}
    \label{fig:impact-panel}
\end{figure}

\subsubsection{Per-LSOA Breakdown}

For multi-placement scenarios, a collapsible section groups placements by their LSOA code (carried as metadata on each placement) and reports per-group summaries including charger count, aggregated energy, carbon, peak demand, and borough name. This enables planners to assess spatial concentration and identify areas receiving disproportionate or insufficient investment.

%--------------------------------------------------------------
\subsection{Scenario Contextualisation via DFES Projections}
\label{sec:dfes}

A key analytical contribution is the integration of scenario-level impact estimates with UK Power Networks' Distribution Future Energy Scenarios (DFES)~\cite{ukpn-dfes}, transforming the DFES timeseries from a passive reference panel into a coordinated supply-versus-demand visualisation.

\subsubsection{Data and Geographic Scope}

DFES projections are retrieved at runtime from the UKPN Open Data portal via a paginated REST API. Each record associates a year, a scenario pathway (\emph{Leading the Way}, \emph{Consumer Transformation}, \emph{System Transformation}, \emph{Falling Short}), a geographic identifier, and projected adoption counts for four technology metrics: electric cars, electric vans, heat pumps, and domestic batteries. The panel supports three levels of geographic scoping---Network, Borough, and LSOA---selectable via a toggle control. When the user clicks an H3 cell on the map, the LSOA code and borough name propagate to the DFES panel, which auto-selects the appropriate scope.

\subsubsection{Supply--Demand Overlay}

The scenario's aggregated impact KPIs are converted into DFES-commensurate units. For example, annual energy delivered is converted to the number of EVs supportable ($V_{\text{EV}} = E_{\text{annual}} / 2{,}500~\text{kWh}$), and peak demand to equivalent domestic battery installations ($V_{\text{BAT}} = P_{\text{peak}} / 4~\text{kW}$). The converted value is rendered as a horizontal dashed line across the DFES timeseries chart.

Shaded fill areas between the supply line and each demand pathway curve communicate the supply--demand relationship (Figure~\ref{fig:dfes-overlay}):

\begin{itemize}[nosep]
    \item \textbf{Green fill} (below the supply line): projected demand is within the scenario's supply capacity.
    \item \textbf{Red fill} (above the supply line): projected demand exceeds the supply capacity, indicating unmet demand.
\end{itemize}

\subsubsection{Year of Constriction}

For each DFES pathway, the platform identifies the \emph{year of constriction}---the first year in which the projected demand curve crosses the supply line from below, computed via linear interpolation:

\begin{equation}
\label{eq:constriction}
t^{*} = t_i + \frac{V - y_i}{y_{i+1} - y_i} \cdot (t_{i+1} - t_i)
\end{equation}

\noindent This is encoded as a filled circle at the intersection point with a vertical drop-line to the $x$-axis, providing a direct temporal answer to the planning question: ``Under which pathway, and in which year, will this deployment become insufficient?''

\begin{figure}[htbp]
    \centering
    % PLACEHOLDER: Screenshot of DFES timeseries with supply-demand overlay
    \fbox{\parbox{0.95\columnwidth}{\centering\vspace{3cm}\textbf{[Figure: DFES Overlay]}\\\small DFES timeseries showing four demand pathways (coloured lines) with the scenario supply line (dashed horizontal). Green shading indicates sufficient capacity; red shading indicates unmet demand. Year-of-constriction markers appear at pathway--supply intersections.\vspace{0.5cm}}}
    \caption{The DFES timeseries panel with supply--demand overlay. The horizontal dashed line represents the scenario's supply capacity (converted to DFES metric units). Shaded regions encode the temporal supply--demand relationship. Intersection markers indicate the year of constriction for each demand pathway.}
    \label{fig:dfes-overlay}
\end{figure}

%--------------------------------------------------------------
\subsection{Scenario Management}
\label{sec:scenario-mgmt}

The Scenario Manager enables users to save the complete analytical state---MCDA weights, method, criterion polarities, AHP comparisons, and all placements---as a named scenario. Saved scenarios can be restored, compared side-by-side via the radar chart, imported/exported as JSON files, and toggled between visible, highlighted, and hidden display modes. This supports iterative planning workflows where a planner evaluates multiple deployment strategies and selects the preferred alternative based on multi-dimensional performance comparison.

%--------------------------------------------------------------
\subsection{Coordinated Linking Model}
\label{sec:linking}

Interactions propagate across views through a shared reactive store:

\begin{itemize}[nosep]
    \item \textbf{Map $\rightarrow$ DFES}: clicking an H3 cell sets the selected LSOA and borough; the DFES panel re-fetches data scoped to the selected geography.
    \item \textbf{Map $\rightarrow$ Impact}: placing a charger triggers impact recalculation; KPI cards and the radar chart update within the same render cycle.
    \item \textbf{Weight Lab $\rightarrow$ Map}: adjusting a criterion weight re-runs the DuckDB MCDA query; the choropleth and glyph overlays update after a 200\,ms debounce to maintain responsive interaction.
    \item \textbf{DFES $\rightarrow$ Map}: selecting an LSOA in the DFES panel updates the map's selected geography and adjusts the scope accordingly.
\end{itemize}

\noindent This architecture supports \emph{loose coupling}: each panel subscribes only to the state slices it needs, minimising unnecessary re-renders during high-frequency interactions such as weight dragging.

% 5_use_case.tex

\section{Use Case: EVCP Planning in Greater London}
\label{sec:use-case}

This section presents a walkthrough of the platform through a representative planning scenario, illustrating how a transport planner might use the system to evaluate and compare alternative EVCP deployment strategies. The walkthrough follows a complete analytical cycle: weight configuration, spatial exploration, placement, impact assessment, temporal contextualisation, and scenario comparison.

%--------------------------------------------------------------
\subsection{Problem Context}
\label{sec:problem-context}

A transport planning officer at a London borough is tasked with identifying suitable locations for new rapid charging infrastructure. The planner must balance several competing objectives: locating chargers near high-demand areas (population density, traffic volume), ensuring equitable coverage in deprived communities, avoiding locations where the electricity grid is already constrained, and complementing the existing public charging network rather than duplicating it. The planner must also estimate the downstream impact of proposed deployments to justify investment to council decision-makers.

%--------------------------------------------------------------
\subsection{Step 1: Configuring Criterion Weights}
\label{sec:uc-weights}

The planner begins by examining the nine criteria in the Weight Laboratory (Figure~\ref{fig:dashboard-layout}). Initially, all criteria carry equal weights ($w_i = 1/9$). The planner adjusts the weights using the parallel coordinate interface, increasing the importance of population density ($C_1$), car ownership ($C_2$), and distance to nearest existing EVCP ($C_9$) while reducing the weight on employment accessibility ($C_4$). Each drag interaction triggers a re-computation of the MCDA scores across all 118,000 H3 cells, and the choropleth map updates within approximately 200\,ms.

To validate the subjective weight assignment, the planner switches to the AHP pairwise comparison panel (Figure~\ref{fig:ahp-interface}). Using the slider controls, the planner expresses that population density is ``moderately more important'' (scale value 3) than car ownership, and ``strongly more important'' (scale value 5) than employment accessibility. After completing the pairwise comparisons, the system reports a Consistency Ratio of $\text{CR} = 0.06 < 0.10$, indicating acceptable consistency. The planner applies the AHP-derived weights, which replace the manually assigned values and immediately update the spatial view.

%--------------------------------------------------------------
\subsection{Step 2: Exploring Spatial Suitability}
\label{sec:uc-spatial}

With the weights configured, the planner examines the suitability map. At the city-wide zoom level, the choropleth displays H3 cells at resolution~7 ($\sim$2,000 hexagons), providing an overview of spatial suitability patterns across Greater London. Areas of high suitability (dark colours on the \texttt{viridis} scale) are concentrated in inner London boroughs with high population density and traffic volume.

The planner activates the glyph overlay at resolution~8 and selects the rose glyph type to examine the multi-criteria profiles of high-scoring areas (Figure~\ref{fig:glyph-map}). Two clusters of high-scoring cells emerge: one in the Southwark area, where all criteria contribute strongly, and another in Hackney, where high population density and deprivation scores dominate but grid headroom is relatively constrained (indicated by a shorter petal for $C_7$). The planner clicks on a glyph in Southwark to pin it as a reference, then hovers over the Hackney glyph to compare profiles. The dashed overlay reveals that the Hackney cell has lower grid headroom and slightly less existing EVCP coverage.

The planner also switches between MCDA methods to assess ranking robustness. Under WSM, the two areas rank similarly. Under WPM, the Hackney cell drops in rank because its low grid headroom value ($\bar{x}_{h,7} = 0.22$) is penalised multiplicatively. Under TOPSIS, the Southwark cell ranks higher due to its closer proximity to the ideal solution across all criteria. This methodological sensitivity analysis, performed in seconds through the method dropdown, informs the planner that the Southwark location is robust across analytical methods while the Hackney location is more sensitive to the treatment of its weak criterion.

%--------------------------------------------------------------
\subsection{Step 3: Placing Chargers and Evaluating Impact}
\label{sec:uc-placement}

The planner activates Simulation Mode and zooms to the Southwark area at resolution~10 to view individual cells. Clicking a cell near a major intersection opens the placement configuration popover, which displays the reverse-geocoded address and the cell's raw criterion values. The planner selects ``DC Rapid'' (50\,kW) as the charger type and sets the count to 4. The estimated installation cost (\pounds160,000) appears immediately.

Upon confirming the placement, the Impact Panel on the right updates with eight KPI cards (Figure~\ref{fig:impact-panel}):

\begin{itemize}[nosep]
    \item Energy delivered: 456,250 kWh/year
    \item Carbon saved: 312 tCO\textsubscript{2}/year
    \item Installation cost: \pounds160,000
    \item Annual revenue: \pounds159,688/year
    \item Peak demand: 140 kW
    \item Headroom impact: 12.3\%
    \item Population served: 1,842 persons
    \item Utilisation: 78\%
\end{itemize}

The planner then places a second deployment of 2 AC Fast chargers (22\,kW) in Hackney. The Impact Panel aggregates the metrics across both placements: total energy delivered rises to 714,570 kWh/year, and total carbon saved to 489 tCO\textsubscript{2}/year. The per-LSOA breakdown shows the contribution of each location. The headroom impact for the Hackney placement is 28.4\%, higher than Southwark's 12.3\%, confirming the grid constraint concern identified during the glyph comparison.

%--------------------------------------------------------------
\subsection{Step 4: Temporal Contextualisation via DFES}
\label{sec:uc-dfes}

The planner examines the DFES timeseries panel below the map (Figure~\ref{fig:dfes-overlay}), which auto-scopes to the Southwark borough based on the selected cell. The electric cars metric shows four demand pathways rising from approximately 2,000 vehicles in 2025 to between 15,000 and 45,000 by 2050. The scenario supply line---computed from the total annual energy delivered, converted to the number of EVs supportable ($V_{\text{EV}} = 714{,}570 / 2{,}500 \approx 286$ vehicles)---appears as a horizontal dashed line.

The green/red shading reveals that the proposed deployment covers the projected demand under all pathways until approximately 2031 (``Leading the Way'' pathway), 2033 (``Consumer Transformation''), 2036 (``System Transformation''), and 2042 (``Falling Short''). The year-of-constriction markers with vertical drop-lines make these transition points immediately legible.

The planner widens the scope to Network level to see the aggregate impact against London-wide demand, observing that the two-site deployment is a small fraction of the city's projected need. Switching to the Borough scope for Hackney reveals tighter headroom: the year of constriction under ``Leading the Way'' occurs in 2029, two years earlier than Southwark, reflecting Hackney's faster projected EV uptake and the smaller deployment size.

%--------------------------------------------------------------
\subsection{Step 5: Saving and Comparing Scenarios}
\label{sec:uc-compare}

Satisfied with the analysis, the planner saves the current scenario as ``Scenario A: Rapid + Fast Mix.'' The planner then creates a second scenario by removing the Hackney placement and adding a larger deployment of 6 DC Ultra-rapid chargers (150\,kW) at a single location near Canary Wharf. This scenario is saved as ``Scenario B: Concentrated Ultra-rapid.''

Both scenarios are selected for comparison in the Scenario Manager. The radar chart in the Impact Panel (Figure~\ref{fig:impact-panel}) normalises the six comparison axes across both scenarios, revealing their distinct profiles:

\begin{itemize}[nosep]
    \item Scenario A (distributed, mixed) scores higher on population coverage and has lower headroom impact but delivers less total energy.
    \item Scenario B (concentrated, ultra-rapid) delivers substantially more energy and carbon savings but concentrates grid stress at a single substation and serves fewer residents directly.
\end{itemize}

\noindent The radar chart enables the planner to assess these trade-offs at a glance. The planner exports both scenarios as JSON files for discussion at the next borough planning meeting, where stakeholders can import them into the platform and explore the assumptions interactively.

%--------------------------------------------------------------
\subsection{Summary}
\label{sec:uc-summary}

This walkthrough demonstrates how the platform supports a complete planning workflow: from criterion weight configuration and spatial exploration through placement simulation, impact estimation, temporal stress testing, and multi-scenario comparison. The interactive nature of the system enables rapid iteration---the planner completed the entire analysis in a single session without batch processing or server-side computation---while the coordinated views ensure that spatial, analytical, and temporal perspectives remain synchronised throughout the exploration.

% 6_discussion.tex

\section{Discussion}
\label{sec:discussion}

%--------------------------------------------------------------
\subsection{Contributions to Visual Analytics for Spatial Decision Support}
\label{sec:disc-va}

The platform demonstrates a visual analytics design pattern that extends the conventional GIS-MCDA workflow in two directions. First, it closes the loop between site selection and impact evaluation: rather than treating suitability scoring as the terminal analytical step, the system enables planners to test the consequences of acting on suitability rankings by simulating charger deployments and observing their estimated effects. Second, it introduces a temporal dimension through the DFES overlay, which transforms static impact estimates into forward-looking assessments by revealing when a deployment becomes insufficient under different demand trajectories. Together, these extensions shift the analytical paradigm from ``where is the best location?'' to ``what happens if we build here, and for how long will it be sufficient?''---a question more closely aligned with the iterative, uncertain nature of real infrastructure planning.

The coordinated multiple-view architecture is central to supporting this extended workflow. The tight coupling between the Weight Laboratory, Spatial View, Impact Panel, and DFES panel enables rapid analytical iteration: a weight change propagates through the MCDA scoring, updates the map, and (if placements exist) cascades into the impact and temporal views within a single interaction cycle. This immediacy supports the exploratory, hypothesis-driven mode of inquiry that visual analytics is designed to facilitate~\cite{keim2008visual}.

%--------------------------------------------------------------
\subsection{Role of Multi-Criteria Method Comparison}
\label{sec:disc-methods}

Offering WSM, WPM, and TOPSIS within a single interface serves as an implicit sensitivity analysis. Rankings that remain stable across methods are robust to aggregation assumptions; rankings that shift---as illustrated in the use case where WPM penalised a location with low grid headroom more severely than WSM---reveal criteria that act as structural constraints rather than compensable factors. This capability has practical value for planning: it helps distinguish ``universally good'' locations from those whose attractiveness depends on how trade-offs are modelled.

The integration of AHP alongside direct weight manipulation similarly broadens the analytical space. AHP provides a structured, theoretically grounded procedure for weight derivation that can be documented and audited, while direct manipulation enables rapid exploratory adjustment. The ability to switch between these modes---and to observe the spatial consequences immediately---supports both the rigour expected in formal planning processes and the fluidity required for exploratory analysis.

%--------------------------------------------------------------
\subsection{Client-Side Architecture}
\label{sec:disc-architecture}

Running the entire analytical pipeline client-side in DuckDB-WASM offers practical advantages. The system can be deployed as a static web application without backend infrastructure, reducing operational cost and simplifying distribution to planning teams. Interactive performance is adequate: the MCDA scoring query over 118,000 rows completes in under 200\,ms on modern hardware, and the H3 hierarchical aggregation for multi-scale visualisation further reduces rendering load.

The use of H3 hexagonal indexing as the spatial unit provides several benefits beyond those described in Section~\ref{sec:spatial-data}. The uniform cell geometry eliminates the area distortion present in administrative boundaries (e.g., LSOAs), enabling consistent spatial aggregation and comparison. The hierarchical structure supports zoom-dependent rendering without pre-computed tile pyramids. The six-neighbour topology supports the 1-ring population estimation in the impact model (Section~\ref{sec:demand}), providing a natural definition of ``service area'' for a charger deployment.

However, the client-side architecture also introduces constraints. The data must be pre-processed into Parquet files and pre-normalised before deployment, limiting the system's ability to incorporate new data sources without a rebuild pipeline. Memory consumption scales with dataset size; the current 118,000-cell dataset consumes approximately 50\,MB of browser memory, which is manageable on desktop devices but may be restrictive on mobile platforms. Future work could explore progressive loading strategies or server-assisted caching for larger study areas.

%--------------------------------------------------------------
\subsection{Limitations of the Impact Model}
\label{sec:disc-limitations}

The impact estimation model (Section~\ref{sec:impact}) is designed for interactive scenario exploration rather than precise forecasting. Several simplifications constrain its applicability:

\begin{itemize}[nosep]
    \item \textbf{Static assumptions}: EV adoption rate ($\alpha_{\text{EV}} = 0.15$), daily charging probability ($\alpha_{\text{charge}} = 0.3$), and grid emission factor ($\text{EF}_{\text{grid}} = 0.233$) are fixed constants. In practice, these vary over time and by location, and the decarbonisation of the electricity grid will reduce the carbon benefit of EV charging.
    \item \textbf{No network effects}: the model treats each placement independently, without accounting for how multiple nearby deployments might compete for demand or collectively affect the local distribution network.
    \item \textbf{Headroom estimation}: substation capacity is estimated from a normalised headroom score rather than from engineering-grade network data. This provides a useful directional signal but should not be used for grid connection assessments.
    \item \textbf{Demand model}: the chain of demographic ratios in the demand equation (Eq.~\ref{eq:demand}) is a rough approximation; actual charging demand depends on trip patterns, dwell times, and driver behaviour that are not captured.
\end{itemize}

\noindent These limitations are acceptable in the context of a visual analytics tool whose purpose is to enable comparative scenario exploration rather than to produce definitive engineering assessments. The transparency of the model---with clearly stated assumptions and visible equations---allows domain experts to evaluate its suitability for their planning context.

%--------------------------------------------------------------
\subsection{Generalisability}
\label{sec:disc-generalisability}

Although demonstrated for EVCP siting in London, the design pattern generalises to other spatial planning problems that share three characteristics: (1) a multi-criteria evaluation of candidate locations, (2) an estimable impact of interventions at selected locations, and (3) a temporal demand projection against which the intervention can be stress-tested. Examples include solar panel deployment, public transit stop placement, urban greenspace allocation, and healthcare facility siting. The modular architecture---where the MCDA engine, impact model, and temporal overlay are independent components linked through a shared spatial key (the H3 cell)---supports adaptation to new domains by replacing the criteria layers, impact equations, and temporal reference data.

% 7_conclusion.tex

\section{Conclusion}
\label{sec:conclusion}

We have presented a visual analytics platform that integrates multi-criteria spatial suitability analysis, what-if impact simulation, and temporal scenario contextualisation for electric vehicle charging infrastructure planning. The system enables planners to interactively adjust criterion weights---through direct manipulation or structured AHP pairwise comparison---observe the resulting suitability rankings across approximately 118,000 H3 hexagonal cells, place candidate charger deployments, evaluate their estimated impact on energy delivery, carbon reduction, grid stress, and population coverage, and contextualise these impacts against Distribution Future Energy Scenarios demand projections.

The platform introduces a design pattern that extends the conventional GIS-MCDA workflow beyond site selection into impact evaluation and temporal stress testing. By overlaying scenario supply on DFES demand trajectories, the system answers a question that static suitability maps cannot: under different demand growth pathways, in which year will a proposed deployment become insufficient? This temporal dimension transforms planning from a one-time ranking exercise into an iterative process of scenario exploration and comparison.

The client-side architecture, built on DuckDB-WASM for analytical computation and H3 for spatial indexing, demonstrates that complex multi-criteria spatial analysis can be performed interactively in the browser without server infrastructure. This lowers the deployment barrier for planning teams and enables the system to function as a self-contained analytical tool that can be distributed alongside the datasets it analyses.

We demonstrated the platform through a case study of EVCP planning in Greater London, using nine heterogeneous open datasets. The use case walkthrough illustrated a complete planning cycle---weight configuration, spatial exploration via choropleth and glyph maps, multi-method sensitivity analysis, charger placement, impact estimation, DFES contextualisation, and multi-scenario comparison via radar charts. The results suggest that the interactive integration of these analytical stages enables a more transparent and exploratory approach to infrastructure planning than batch-oriented GIS-MCDA workflows.

Future work will address several directions. First, a formal user study with transport planners will evaluate the platform's usability and the decision quality it supports compared to existing workflows. Second, the impact model can be enhanced with time-varying assumptions (e.g., evolving grid emission factors and EV adoption rates) to improve the fidelity of temporal projections. Third, the modular architecture could be adapted to other spatial planning domains---such as renewable energy siting, healthcare facility placement, or urban greenspace allocation---by substituting the criteria layers, impact equations, and temporal reference data. Finally, collaborative features such as shared scenario libraries and annotation tools could extend the platform's utility for group decision-making processes~\cite{jankowski2001gis}.


\bibliographystyle{IEEEtran}
\bibliography{references}

\end{document}
